27.

(a) The matrix of $\phi^{-1}$ with respect to the basis $\{m_\lambda\}$ is the matrix $M$ of Section 7.4. By a simple result of linear algebra, $M$ and $M^{-1}$ have the same Jordan block sizes. By Theorem 7.4.4 the matrix $M$ is upper triangular with 1's on the main diagonal. Hence the size of the largest Jordan block of $\phi$ is the least integer $m$ for which $(M-I)^m=0$. Again by Theorem 7.4.4 we have $M_{\lambda\mu}=0$ unless $\mu\le\lambda'$. By Proposition 7.4.1 and the Gale-Ryser theorem (a basic combinatorial result which everyone should know) the converse is true, i.e., $M_{\lambda\mu}>0$ if $\mu\le\lambda'$. Since all entries of $M$ are nonnegative, there is no cancellation in the computation of powers of $M-I$. It follows that $m$ is the size of the longest chain in the poset $\mathrm{Par}(n)$, ordered by dominance.

Now use Exercise 7.2(f).

(This exercise gives the \emph{length} of the longest chain, so you need to add 1 to get the size.) For some further work on chains in $\mathrm{Par}(n)$, see E. Early, \textit{Discrete Math.} \textbf{313} (2013), 2168--2177.
